\chapter{Внутренний шум, модульное время и предел столкновений}
\label{ch:v8-intrinsic-noise-clock-dilation-gate}

\section{Безразмерный поток уже существует}

Для выведенного генератора можно без дополнительных структур записать
\begin{equation}
 T_u=\exp(u\mathcal L_{q\ell}),
 \qquad u\geq0.
 \label{eq:v8-dimensionless-lindblad-flow}
\end{equation}
На полной концевой алгебре межсекторный генератор имеет $46$ нулевых направлений,
наименьшую ненулевую величину $1/2$ и максимальную величину $8$ в принятой
нормировке. Добавление связывающей и калибровочных форм ранее сокращает неподвижную
алгебру до скаляра.

Но положительное масштабирование
\begin{equation}
 \mathcal L_{q\ell}\longmapsto\kappa\mathcal L_{q\ell}
 \label{eq:v8-rate-rescaling}
\end{equation}
не меняет Краус-пространство, калибровочная ковариантность или неподвижную
алгебру, а все скорости умножает на $\kappa$. Поэтому $u$ является
каноническим безразмерным параметром полугруппы, но ещё не физическим
временем.

\section{Почему собственное модульное время не создаёт шум}

Для стационарного следового состояния $\rho_*=I_{21}/21$ модульный
гамильтониан $-\log\rho_*$ скалярен, и
\begin{equation}
 i[-\log\rho_*,X]=0.
\end{equation}
Более общий верный элемент уже выбранной центральной алгебры имеет вид
\begin{equation}
 \rho_{q\ell}=aP_q+bP_\ell,
 \qquad a,b>0.
\end{equation}
Его модульный поток может приписывать фазы межсекторным операторам, но
точно фиксирует оба проекторных населения:
\begin{equation}
 [\log\rho_{q\ell},P_q]
 = [\log\rho_{q\ell},P_\ell]=0.
\end{equation}
В то же время диссипативный генератор действует на $P_q$ ненулево. Значит,
собственный обратимый модульный поток состояния не совпадает с требуемым
необратимым Краус-процессом. Это повторяет, теперь на новом операторном
носителе, границу Тома VI.

\section{Столкновительный предел}

Одношаговый канал предыдущей главы действительно аппроксимирует поток,
если каждый раз предоставляется свежая тринадцатимерная вспомогательная система и
\begin{equation}
 p=\frac{u}{n},\qquad n\longrightarrow\infty.
 \label{eq:v8-fresh-ancilla-scaling}
\end{equation}
Для $u=0.1$ машинная ошибка относительно
$\exp(u\mathcal L_{q\ell})$ монотонно падает от $2.1\cdot10^{-1}$ при
одном столкновении до величины меньше $10^{-3}$ при тысяче столкновений.
Начальная ошибка равна $1.008$, а не считается малым возмущением; сходимость
возникает именно в пределе многих слабых столкновений.

Это конструктивный результат: непрерывный квантовый шум получается как
предел уже выведенных конечных каналов. Однако сама бесконечная цепь свежих
состояний вспомогательных систем, длительность одного такта и коэффициент $\kappa$ пока
не являются следствиями конечного родителя.

\begin{theorem}[Безразмерное шумовое время]
Геометрия межсекторных стрелок канонически определяет полугруппу Линдблада с
безразмерным параметром и допускает её столкновительную аппроксимацию на
повторениях минимального Стайнспринг-носителя. Но ни собственный модульный
поток стационарного состояния, ни ранее построенные четырёхтактные часы не
порождают требуемую диссипацию и не фиксируют физическую скорость.
\end{theorem}

\begin{nogocriterion}[Запрет калибровки времени собственной щелью]
Нельзя положить наименьшую скорость равной единице и считать секунду
выведенной. Такая операция лишь определяет единицу безразмерного параметра
$u$; преобразование к физическому времени требует независимо полученного
$\kappa$ или часового взаимодействия.
\end{nogocriterion}

\begin{protocol}
\begin{enumerate}
 \item каноническая безразмерная полугруппа: \textbf{получена};
 \item пространство кратностей шума: \textbf{получено};
 \item положительное масштабирование сохраняет структуру: \textbf{да};
 \item физическая скорость $\kappa$: \textbf{не выведена};
 \item следовое модульное время: \textbf{тривиально};
 \item центральное верное состояние переносит населения: \textbf{нет};
 \item предел свежих вспомогательных систем: \textbf{сходится};
 \item источник бесконечной цепи вспомогательных систем: \textbf{не выведен};
 \item четырёхтактные проектные часы: \textbf{не являются источником шума};
 \item автономное физическое время: \textbf{не получено};
 \item следующий тест: \textbf{детального баланса селектор относительных скоростей}.
\end{enumerate}
\end{protocol}

Машинный сертификат сохранён в
\path{s2t_v8_intrinsic_noise_clock_dilation_gate_results.json}.

\begin{thebibliography}{9}
\bibitem{v8-clock-connes-rovelli}
A.~Connes, C.~Rovelli,
\emph{Von Neumann algebra automorphisms and time--thermodynamics relation
in generally covariant quantum theories},
Classical and Quantum Gravity \textbf{11} (1994), 2899--2918;
arXiv:gr-qc/9406019.
\bibitem{v8-clock-attal-joye}
S.~Attal, A.~Joye,
\emph{Weak coupling and continuous limits for repeated quantum
interactions}, Journal of Statistical Physics \textbf{126} (2007),
1241--1283; arXiv:math-ph/0501012.
\bibitem{v8-clock-attal-pautrat}
S.~Attal, Y.~Pautrat,
\emph{From repeated to continuous quantum interactions},
Annales Henri Poincar\'e \textbf{7} (2006), 59--104;
arXiv:math-ph/0311002.
\end{thebibliography}